\section{Two rows}\label{sec:strip}

On $2$ rows every rectangle is axis-aligned, since a tilted rectangle meets at
least $3$ rows by Lemma~\ref{lem:leftmost}. A rectangle of the $2\times m$ grid
is therefore given by its $2$ columns.

\begin{theorem}[Strip Theorem]\label{thm:strip}
Let $m\ge3$.
\begin{enumerate}[(a)]
\item A blackout of the $2\times m$ grid is valid if and only if at most $1$ of
its columns is full.
\item We have $f(2,m)=m+1$, and the optima are the blackouts with $1$ full column
and $1$ black point in every other column.
\item We have $a_2(m)=m\,2^{m-1}$. So row $2$ of \seqnum{A399580} is
\seqnum{A001787} from $m=3$ on.
\end{enumerate}
For $m=2$ there is $1$ rectangle, so $f(2,2)=4$ and $a_2(2)=1$.
\end{theorem}

\begin{proof}
(a) Take two distinct rectangles, with column sets $\{p,q\}$ and $\{p',q'\}$.
Their corner sets differ exactly in the points of the columns that belong to
one of these sets and not to the other. There are $2$ or $4$ such columns.
Conversely, for any two columns $p\ne q$ choose a third column $w$. Here we
use $m\ge3$. The rectangles on $\{w,p\}$ and $\{w,q\}$ differ exactly in
the points of columns $p$ and $q$. By Lemma~\ref{lem:diff}, a blackout is
valid if and only if no $2$ columns are full.

(b, c) A valid blackout keeps a white point in all but at most $1$ column, so
$|S|\le2m-(m-1)=m+1$. Equality requires $1$ full column and exactly $1$ black
point in each other column. There are $m$ choices for the full column and $2$
choices in each of the other $m-1$ columns.
\end{proof}

The $2\times m$ grid has $4$ symmetries, generated by the row reflection and
the column reversal. We count the optima up to these.

\begin{proposition}[orbits and kings]\label{prop:kings}
For $m\ge3$ the number of optima of the $2\times m$ grid up to symmetry is
\[
m\,2^{m-3}+[m\text{ odd}]\,2^{(m-3)/2}.
\]
It equals \seqnum{A322284}$(m-1)$, the number of ways to place $m-1$
nonattacking kings on a $2\times(2m-2)$ board up to symmetry.
\end{proposition}

\begin{proof}
We use Burnside's lemma. The identity fixes all $m2^{m-1}$ optima, and the row
reflection fixes none, since it moves every singleton. The column reversal and
the half turn can fix an optimum only when $m$ is odd and the full column is
the middle one. Each of them then fixes $2^{(m-1)/2}$ optima, since the
singleton columns come in mirror pairs and one column of each pair determines
the other. Dividing the sum by $4$ gives the formula.

For the kings, put $n=m-1$. Two kings on a board with $2$ rows attack each other
exactly when their columns are equal or adjacent. So a placement of $n$
nonattacking kings on the $2\times2n$ board consists of columns
$c_1<\cdots<c_n$ with $c_{j+1}\ge c_j+2$, and a row $r_j$ for each king. The
numbers $c_1-1$, $c_{j+1}-c_j-2$ for $1\le j<n$, and $2n-c_n$ are nonnegative and
sum to $1$. Hence exactly one of these $n+1$ numbers is $1$. If it is the $e$th,
then $c_j=2j-1+[j\ge e]$. Send this placement to the optimum with full column
$e$ and singleton rows $r_1,\ldots,r_n$ from left to right. This is a bijection.
It commutes with the row reflection. The column reversal of the board sends
$c$ to $2n+1-c$ and reverses the order of the rows $r_j$. Since
$2n+1-c_{n+1-j}=2j-1+[j\ge n+2-e]$, it sends $e$ to $n+2-e=m+1-e$, as the column
reversal of the grid does. Hence the orbits correspond.
\end{proof}
